\noindent (c) 哪几对线性组合有零协方差?

\medskip
\noindent 2.42\quad 重复练习 2.41, 但
\[
\boldsymbol{\Sigma}_{\mathbf{X}}=\begin{bmatrix}3&1&1&1\\1&3&1&1\\1&1&3&1\\1&1&1&3\end{bmatrix}
\]

\vspace{18pt}
\begin{center}
\phantomsection\label{bib:2}
{\large\bfseries 参考文献}
\end{center}
\vspace{4pt}

\begin{enumerate}[label=\arabic*.,leftmargin=*]
\item Bellman, R. \emph{Introduction to Matrix Analysis} (2nd ed.) Philadelphia: Soc for Industrial \& Applied Math (SIAM), 1997.
\item Eckart, C., and G. Young. ``The Approximation of One Matrix by Another of Lower Rank.'' \emph{Psychometrika}, \textbf{1} (1936), 211--218.
\item Graybill, F. A. \emph{Introduction to Matrices with Applications in Statistics}. Belmont, CA: Wadsworth, 1969.
\item Halmos, P. R. \emph{Finite-Dimensional Vector Spaces}. New York: Springer-Verlag, 1993.
\item Johnson, R. A., and G. K. Bhattacharyya. \emph{Statistics: Principles and Methods} (5th ed.) New York: John Wiley, 2005.
\item Noble, B., and J. W. Daniel. \emph{Applied Linear Algebra} (3rd ed.). Englewood Cliffs, NJ: Prentice Hall, 1988.
\end{enumerate}
